\section{Observation model, identification, and decision criteria}

Generic constants \(c\) and \(C\) denote positive numerical constants whose values may change across occurrences. Absolute values denote scalar magnitudes, and \(a\asymp b\) means that \(a\) is bounded above and below by positive numerical multiples of \(b\). For sample size \leanref{sym:n}{\ensuremath{n}}, record indices run from \(1\) to \(n\); sums over empty index sets equal zero. Expectations and probabilities with procedure subscripts integrate the procedure's randomization as well as the observed sample.

\subsection{Causal records and observational sampling}

The full causal record is \((X,A,(Y(0),Y(1)),Y)\), where \(X\in[0,1]\) is the scalar covariate, \(A\in\{0,1\}\) is treatment, and the potential and observed outcomes are binary. Its observed marginal retains \((X,A,Y)\), and an ordered sample consists of \(n\) independent records from that marginal. The anchored definitions immediately below record the coordinate, marginal-law, and dataset conventions used by the statistical class and decision problems.

The model uses a common public regularity radius \leanref{sym:L}{\ensuremath{L}} for the nuisance functions and the treatment-effect contrast. Its numerical value is fixed throughout the analysis.

\begin{definitionv}{synth_4}[Fixed public regularity radius]
The public regularity radius \(L\) is the real constant
\[
L = 3.
\]
\end{definitionv}

The causal record contains the covariate \leanref{sym:X}{\ensuremath{X}}, treatment indicator \leanref{sym:A}{\ensuremath{A}}, potential-outcome pair \leanref{sym:Ypot}{\ensuremath{Y^{\mathrm{pot}}}}, and observed outcome \leanref{sym:Y}{\ensuremath{Y}}. The following coordinate conventions distinguish the observed variables from the potential outcomes.

\begin{definitionv}{synth_9}[Observed records and ordered datasets]
The observed-record space is
\[
\mathcal O=[0,1]\times\{0,1\}^2,
\]
equipped with its Borel sigma-algebra, with a record written as \(z=(x,a,y)\). For a nonnegative integer \(n\), define \(\mathcal O^n\) to be its ordered \(n\)-fold product, equipped with the product sigma-algebra. A dataset is an element
\[
D=(z_1,\ldots,z_n)\in\mathcal O^n.
\]
For a random dataset \(D\), write \(\mathcal L(D)\) for its probability law on \(\mathcal O^n\); explicitly,
\[
\mathcal L(D)(E)=\Pr\{D\in E\}
\]
for every measurable \(E\subseteq\mathcal O^n\). For \(n=0\), the product consists of the empty dataset.
\end{definitionv}

\begin{definitionv}{synth_3}[Observed binary outcome coordinate]
The observed binary outcome \(Y\) is the final coordinate of a full causal record, whose coordinates are, in order, a covariate, a binary treatment indicator, an ordered pair of binary potential outcomes, and a binary observed outcome.
\end{definitionv}

\begin{definitionv}{synth_2}[Binary treatment coordinate]
The binary treatment indicator \(A\) maps each full causal record to its treatment coordinate, taking values in \(\{0,1\}\). A full causal record is ordered as covariate, binary treatment, ordered pair of binary potential outcomes, and binary observed outcome.
\end{definitionv}

A causal law \leanref{sym:P}{\ensuremath{P}} specifies the joint distribution of the full record. Its observational marginal \leanref{sym:Pobs}{\ensuremath{P^{\mathrm{obs}}}} specifies the distribution available to an estimator; the covariate marginal \(P_X\) describes the random design.

\begin{definitionv}{synth_1}[Scalar covariate coordinate]
The scalar covariate \(X\) is the first coordinate of the causal record, with values in the closed unit interval:
\[
X \in [0,1].
\]
\end{definitionv}

Write \(Y(A)=(1-A)Y(0)+AY(1)\) for the potential outcome selected by the realized treatment. The observed-record space \leanref{sym:O}{\ensuremath{\mathcal O}} retains the covariate, treatment, and observed outcome, and an ordered dataset \leanref{sym:D}{\ensuremath{D}} collects these records.

\begin{definitionv}{synth_8}[Causal coordinates and marginal laws]
A causal law \(P\) is a Borel probability law of
\[
(X,A,Y^{\mathrm{pot}},Y)
\in[0,1]\times\{0,1\}\times\{0,1\}^2\times\{0,1\},
\]
where \(X\) is the covariate, \(A\) is the treatment indicator, \(Y\) is the observed outcome, and
\[
Y^{\mathrm{pot}}=(Y(0),Y(1))
\]
is the ordered pair of control and treatment potential outcomes. Define the observational marginal and the covariate marginal by
\[
P^{\mathrm{obs}}=\mathcal L_P(X,A,Y),
\qquad
P_X=\mathcal L_P(X).
\]
Here \(\mathcal L_P\) denotes the probability law of the indicated coordinates under \(P\).
\end{definitionv}

Sampling and causal identification rest on the following three conditions. The first specifies how observed records enter the dataset.

\begin{assumptionv}{ass:iid}[Independent and identically distributed observations]
The dataset \(D\) has distribution
\[
\mathcal L(D)=(P^{\mathrm{obs}})^{\otimes n}.
\]
\end{assumptionv}

Independent and identically distributed observational sampling makes the common observed-record law the statistical experiment for estimation and inference \citep{KennedyBalakrishnanRobinsWasserman2024}. The next condition connects each observed outcome to the corresponding potential outcome.

\begin{assumptionv}{ass:consistency}[Consistency of observed outcomes]
The observed outcome \(Y\) equals the potential outcome \(Y(A)\) selected by the realized treatment with probability one:
\[
P\{Y=Y(A)\}=1.
\]
\end{assumptionv}

Consistency assigns the realized outcome to the treatment actually received, giving the observed outcome regression its potential-outcome interpretation \citep{KennedyBalakrishnanRobinsWasserman2024}. Identification further requires that the covariate account for treatment selection in the following sense.

\begin{assumptionv}{ass:exchangeability}[Conditional exchangeability of treatment]
The potential-outcome pair \(Y^{\mathrm{pot}}\) is conditionally independent of the treatment indicator \(A\) given the covariate \(X\) under the causal law \(P\):
\[
Y^{\mathrm{pot}}\perp\!\!\!\perp A\mid X\quad[P].
\]
\end{assumptionv}

Conditional exchangeability permits comparison of the two treatment arms at a common covariate value \citep{KennedyBalakrishnanRobinsWasserman2024}. Together with consistency and the overlap restriction below, it identifies the conditional mean potential-outcome contrast from observational regressions, initially up to covariate-null sets.

\subsection{Selected regressions and the statistical class}

Pointwise estimation requires explicit regression versions. We therefore specify the selected design density \leanref{sym:f}{\ensuremath{f_P}}, propensity regression \leanref{sym:e}{\ensuremath{e_P}}, and control outcome regression \leanref{sym:mu0}{\ensuremath{\mu_{0,P}}} before imposing their regularity conditions.

\begin{definitionv}{synth_10}[Selected design density and observational regressions]
For a causal law \(P\) whose covariate marginal \(P_X\) is absolutely continuous with respect to Lebesgue measure \(dx\) on \([0,1]\), let \(f_P\) be the selected measurable version of its density:
\[
f_P=\frac{dP_X}{dx},
\qquad
P_X(B)=\int_B f_P(x)\,dx
\]
for every Borel set \(B\subseteq[0,1]\). Let \(e_P\) and \(\mu_{0,P}\) be the selected versions of the propensity regression and the control outcome regression, respectively:
\[
e_P(x)=P(A=1\mid X=x),
\qquad
\mu_{0,P}(x)=E_P[Y\mid X=x,A=0],
\qquad x\in[0,1].
\]
These conditional expressions denote the selected regression functions on \([0,1]\). The model's overlap and smoothness conditions apply to these selected versions, and its density bounds apply to \(f_P\) Lebesgue-almost everywhere.
\end{definitionv}

The selected treated outcome regression \leanref{sym:mu1}{\ensuremath{\mu_{1,P}}} completes the observational contrast \leanref{sym:tau}{\ensuremath{\tau_P}}. We use the selected conditional versions throughout and define
\[
\mu_{1,P}(x)=E_P[Y\mid X=x,A=1],
\qquad
\tau_P(x)=\mu_{1,P}(x)-\mu_{0,P}(x),
\qquad
x_0=\frac12,
\qquad
\theta(P)=\tau_P(x_0).
\]
Thus the fixed public evaluation location \leanref{sym:x0}{\ensuremath{x_0}} determines the scalar target \leanref{sym:theta}{\ensuremath{\theta(P)}}.



Binary outcomes place both the contrast and the target in \([-1,1]\). We next organize the random-design, overlap, and smoothness restrictions that determine the statistical class.

\begin{assumptionv}{ass:density}[Bounded covariate density]
The covariate density \(f_P\) satisfies
\[
\frac12\le f_P(x)\le\frac32
\quad\text{for Lebesgue-almost every }x\in[0,1].
\]
\end{assumptionv}

The bounded positive random-design density ensures that covariate neighborhoods receive probability comparable to their Lebesgue length. These fixed numerical bounds instantiate the density condition in \citep[Theorem 1, footnote 6]{WagerAthey2018}. Treatment probabilities are likewise bounded away from their endpoints.

\begin{assumptionv}{ass:overlap}[Uniform treatment overlap]
The conditional treatment probability \(e_P\) satisfies
\[
\frac14\le e_P(x)\le\frac34
\qquad\text{for every }x\in[0,1].
\]
\end{assumptionv}

Uniform treatment overlap ensures that both treatment arms remain represented throughout the covariate domain; the constants give a fixed numerical version of the standard overlap condition \citep{KennedyBalakrishnanRobinsWasserman2024}. The nuisance regressions may vary substantially over short distances, subject to the following common Hölder order.

\begin{assumptionv}{ass:propensity-holder}[Hölder continuity of propensity]
The conditional treatment probability \(e_P\) satisfies the Hölder bound with common smoothness radius \(L\):
\[
|e_P(x)-e_P(x^{\prime})|
\le L|x-x^{\prime}|^{1/10}
\qquad\text{for every }x,x^{\prime}\in[0,1].
\]
\end{assumptionv}

Hölder propensity regularity controls local changes in treatment selection while allowing a rough propensity function of order \(1/10\) \citep{KennedyBalakrishnanRobinsWasserman2024}. The control regression obeys the corresponding restriction.

\begin{assumptionv}{ass:control-holder}[Hölder continuity of control means]
For every \(x,x'\in[0,1]\),
\[
|\mu_{0,P}(x)-\mu_{0,P}(x')|
\le L|x-x'|^{1/10}.
\]
\end{assumptionv}

Hölder control-regression regularity bounds local changes in the baseline outcome mean at the same rough order \citep{KennedyBalakrishnanRobinsWasserman2024}. The treatment-effect contrast has greater regularity than these nuisance functions.

\begin{assumptionv}{ass:contrast-lipschitz}[Lipschitz continuity of contrasts]
For every \(x,x'\in[0,1]\),
\[
|\tau_P(x)-\tau_P(x')|
\le L|x-x'|.
\]
\end{assumptionv}

Lipschitz contrast regularity controls the change in the conditional treatment effect over a covariate neighborhood. It is the Hölder order-one condition under the convention used in \citep{KennedyBalakrishnanRobinsWasserman2024}. Thus the model combines a Lipschitz target function with propensity and control regressions of Hölder order \(1/10\), all within the fixed radius in \cref{obj:synth_4}.

We collect these statistical restrictions into a class of causal laws. The private-release requirement is introduced afterward, once its output-space notation is in place.

\begin{definitionv}{synth_16}[Borel probability laws]
For a space \(\mathsf X\) equipped with its Borel \(\sigma\)-algebra \(\mathcal B(\mathsf X)\), define
\[
\operatorname{Prob}(\mathsf X)
=
\{P:P\text{ is a probability measure on }(\mathsf X,\mathcal B(\mathsf X))\}.
\]
\end{definitionv}

The statistical class includes the identifying conditions and the design and function restrictions just stated.

\begin{definitionv}{def:complete-model}[Causal law class \(\mathcal P\)]
\(\mathcal P\) is the class of causal probability laws
\[
P\in\operatorname{Prob}
\bigl([0,1]\times\{0,1\}\times\{0,1\}^2\times\{0,1\}\bigr)
\]
satisfying all of the following conditions:
\begin{itemize}
\item \textbf{(Consistency.)} \cref{obj:ass:consistency}.
\item \textbf{(Exchangeability.)} \cref{obj:ass:exchangeability}.
\item \textbf{(Covariate density.)} \cref{obj:ass:density}.
\item \textbf{(Overlap.)} \cref{obj:ass:overlap}.
\item \textbf{(Propensity smoothness.)} \cref{obj:ass:propensity-holder}.
\item \textbf{(Control-mean smoothness.)} \cref{obj:ass:control-holder}.
\item \textbf{(Contrast smoothness.)} \cref{obj:ass:contrast-lipschitz}.
\end{itemize}
The function conditions apply to the selected versions of the observational regressions. The covariate density is an unknown measurable function. Both arm means range over their full binary-outcome domain, subject to the stated conditions. Membership in \(\mathcal P\) is determined by these properties of the causal law.
\end{definitionv}

For laws in \cref{obj:def:complete-model}, consistency, exchangeability, and overlap identify the observational contrast with the conditional mean of \(Y(1)-Y(0)\) almost everywhere under the covariate law. Contrast continuity and the positive design density then determine its pointwise interpretation: continuous versions agreeing almost everywhere agree throughout \([0,1]\). Consequently, the target \(\theta(P)\) is the value of the continuous causal contrast at the fixed evaluation location. The argument is given in \cref{obj:lem:point-version}.

\subsection{Private releases and decision criteria}

Privacy constrains a release on the full ordered dataset space in \cref{obj:synth_9}. Its comparison concerns an arbitrary dataset and a comparison dataset \leanref{sym:Dprime}{\ensuremath{D'}}, using replacement Hamming distance \leanref{sym:dHam}{\ensuremath{d_{\mathrm H}(D,D')}}. A randomized release \leanref{sym:M}{\ensuremath{M}} takes values in a standard Borel output space \leanref{sym:B}{\ensuremath{\mathcal B}}; the public privacy budget \leanref{sym:epsilon}{\ensuremath{\epsilon}} bounds the change in its output probabilities.



The privacy requirement applies to every replacement-adjacent pair, including datasets outside any particular sampling law's support.

\begin{assumptionv}{ass:pure-privacy}[Pure replacement differential privacy]
For every pair of datasets \(D,D'\in\mathcal O^n\) with replacement distance \(d_{\mathrm H}(D,D')=1\), and every measurable output event \(E\in\mathcal B(\mathcal B)\),
\[
M(D,E)\le \exp(\epsilon)M(D',E).
\]
\end{assumptionv}

Pure central differential privacy under replacement adjacency bounds the multiplicative change in every output-event probability when one observed record is replaced \citep{AcharyaSunZhang2021}. Requiring this property everywhere makes privacy a condition on the release itself, uniformly across the full dataset space. The admissible class is therefore defined independently of the causal law.

\begin{definitionv}{def:private-kernels}[Private release class \(\mathfrak M_{n,\epsilon}(\mathcal B)\)]
\(\mathfrak M_{n,\epsilon}(\mathcal B)\) is the class of all everywhere-defined Markov kernels
\[
M:\mathcal O^n\rightsquigarrow\mathcal B
\]
satisfying \cref{obj:ass:pure-privacy}. Here \(\mathcal B\) is a standard Borel output space, and \(M(D,E)\) is the conditional probability of an output in the Borel event \(E\) given dataset \(D\).
\end{definitionv}

We use separate decision criteria for scalar estimation and interval inference. For a randomized scalar estimator \leanref{sym:Tgeneric}{\ensuremath{T}}, the minimax absolute-error risk \leanref{sym:R}{\ensuremath{R_{n,\epsilon}}} measures the smallest achievable worst-case expected distance from the pointwise target.

\begin{definitionv}{def:risk-criterion}[Minimax absolute-error risk \(R_{n,\epsilon}\)]
\[
R_{n,\epsilon}
=
\inf_{T\in\mathfrak M_{n,\epsilon}(\mathbb R)}
\sup_{P\in\mathcal P}
E_{(P^{\mathrm{obs}})^{\otimes n},T}
\bigl|T(D)-\theta(P)\bigr|.
\]
Here \(T(D)\) is the scalar output of the randomized estimator, \(\theta(P)\) is the pointwise treatment-effect target, and the expectation integrates both the observed sample with law \((P^{\mathrm{obs}})^{\otimes n}\) and the estimator randomization.
\end{definitionv}

The supremum evaluates accuracy uniformly over the causal class in \cref{obj:def:complete-model}, while the infimum ranges over all private scalar releases in \cref{obj:def:private-kernels}. Interval inference instead evaluates expected length subject to a uniform coverage requirement. To express that criterion for randomized procedures, we first specify the measurable interval output space \leanref{sym:Ispace}{\ensuremath{\mathcal I}}.

\begin{definitionv}{synth_13}[Measurable interval output space]
Define
\[
\mathcal I=\{I\subseteq\mathbb R:I\text{ is a connected Borel set}\}.
\]
Code the empty set by a distinguished symbol \(\varnothing\). Code each nonempty interval by its lower and upper endpoints \(\ell,u\) and endpoint-inclusion flags \(\eta_\ell,\eta_u\in\{0,1\}\), where
\[
\ell\in\mathbb R\cup\{-\infty\},\qquad
u\in\mathbb R\cup\{+\infty\},\qquad \ell\le u.
\]
For \(\ell<u\), the code represents
\[
(\ell,u)\cap\mathbb R
\;\cup\;
\begin{cases}\{\ell\},&\ell\in\mathbb R,\ \eta_\ell=1,\\ \varnothing,&\text{otherwise},\end{cases}
\;\cup\;
\begin{cases}\{u\},&u\in\mathbb R,\ \eta_u=1,\\ \varnothing,&\text{otherwise}.\end{cases}
\]
An infinite endpoint has inclusion flag zero. For \(\ell=u\), require \(\ell=u\in\mathbb R\) and \(\eta_\ell=\eta_u=1\), so the code represents the singleton \(\{\ell\}\). Equip the valid nonempty codes with the Borel sigma-algebra inherited from
\[
(\mathbb R\cup\{-\infty\})\times
(\mathbb R\cup\{+\infty\})\times\{0,1\}^2,
\]
using the usual extended-real Borel structures and discrete flags, and adjoin the empty-set code as an isolated point. Transport this sigma-algebra to \(\mathcal I\) through the coding bijection. This makes \(\mathcal I\) a standard Borel measurable space, including empty, singleton, bounded, and unbounded interval outputs.
\end{definitionv}

This coding accommodates all connected Borel intervals and makes coverage events and interval lengths measurable. For a randomized interval procedure \leanref{sym:Igeneric}{\ensuremath{I}}, the honest expected-length criterion \leanref{sym:H}{\ensuremath{H_{n,\epsilon}}} minimizes maximal expected length among private procedures attaining at least \(90\%\) coverage for every law in the class.

\begin{definitionv}{def:interval-criterion}[Honest expected-length criterion \(H_{n,\epsilon}\)]
The minimax expected length under uniform \(90\%\) coverage is
\[
H_{n,\epsilon}
=
\inf_{\substack{
I\in\mathfrak M_{n,\epsilon}(\mathcal I)\\
\inf_{P\in\mathcal P}
\Pr_{(P^{\mathrm{obs}})^{\otimes n},I}
\{\theta(P)\in I(D)\}\ge 9/10
}}
\ \sup_{P\in\mathcal P}
E_{(P^{\mathrm{obs}})^{\otimes n},I}
\operatorname{Leb}(I(D)).
\]
Here \(\mathcal I\) is the interval output space, \(I(D)\) is the released connected interval, and \(\operatorname{Leb}(I(D))\) is its Lebesgue length. The probability and expectation integrate both the observed sample and the procedure's randomization.
\end{definitionv}

Honesty here means uniform coverage over the entire causal class at the given sample size and privacy budget. Connectedness gives length its usual interpretation as the span of a confidence interval, and both coverage and expected length account for sampling and release randomization. Since the target lies in \([-1,1]\), the constant release of that interval is private, has coverage one, and supplies a finite-length admissible procedure. The two criteria in \cref{obj:def:risk-criterion,obj:def:interval-criterion} thus evaluate estimation accuracy and honest interval precision within the same observational experiment and privacy model.
